\documentclass[12pt]{article}
\usepackage{amsmath,amssymb,amsfonts}
\usepackage[margin=1in]{geometry}
\begin{document}
\begin{center}
{\large\bfseries 6th Irish Mathematical Olympiad\\8 May 1993}
\end{center}
\bigskip\noindent\textbf{Paper 1}\bigskip
\begin{enumerate}
\item The real numbers $\alpha$, $\beta$ satisfy the equations
\begin{align*}
\alpha^3 - 3\alpha^2 + 5\alpha - 17 &= 0,\\
\beta^3 - 3\beta^2 + 5\beta + 11 &= 0.
\end{align*}
Find $\alpha + \beta$.

\item A natural number $n$ is called \textbf{good} if it can be written in a \emph{unique} way simultaneously as the sum $a_1 + a_2 + \ldots + a_k$ and as the product $a_1a_2\ldots a_k$ of some $k \ge 2$ natural numbers $a_1, a_2, \ldots, a_k$. (For example 10 is good because $10 = 5 + 2 + 1 + 1 + 1 = 5 \cdot 2 \cdot 1 \cdot 1 \cdot 1$ and these expressions are unique.) Determine, in terms of prime numbers, which natural numbers are good.

\item The line $l$ is tangent to the circle $S$ at the point $A$; $B$ and $C$ are points on $l$ on opposite sides of $A$ and the other tangents from $B$, $C$ to $S$ intersect at a point $P$. If $B$, $C$ vary along $l$ in such a way that the product $|AB|\cdot|AC|$ is constant, find the locus of $P$.

\item Let $a_0, a_1, \ldots, a_{n-1}$ be real numbers, where $n \ge 1$, and let the polynomial
\[f(x) = x^n + a_{n-1}x^{n-1} + \ldots + a_0\]
be such that $|f(0)| = f(1)$ and each root $\alpha$ of $f$ is real and satisfies $0 < \alpha < 1$. Prove that the product of the roots does not exceed $1/2^n$.

\item Given a complex number $z = x + iy$ ($x, y$ real), we denote by $P(z)$ the corresponding point $(x,y)$ in the plane. Suppose $z_1, z_2, z_3, z_4, z_5, \alpha$ are nonzero complex numbers such that
\begin{enumerate}
\item[(a)] $P(z_1)$, $P(z_2)$, $P(z_3)$, $P(z_4)$, $P(z_5)$ are the vertices of a convex pentagon $\mathbf{Q}$ containing the origin 0 in its interior and
\item[(b)] $P(\alpha z_1)$, $P(\alpha z_2)$, $P(\alpha z_3)$, $P(\alpha z_4)$ and $P(\alpha z_5)$ are all inside $\mathbf{Q}$.
\end{enumerate}
If $\alpha = p + iq$, where $p$ and $q$ are real, prove that $p^2 + q^2 \le 1$ and that
\[p + q\tan(\pi/5) \le 1.\]
\end{enumerate}

\newpage
\begin{center}
{\large\bfseries 6th Irish Mathematical Olympiad\\8 May 1993}
\end{center}
\bigskip\noindent\textbf{Paper 2}\bigskip
\begin{enumerate}
\setcounter{enumi}{5}
\item Given five points $P_1$, $P_2$, $P_3$, $P_4$, $P_5$ in the plane having integer coordinates, prove that there is at least one pair $(P_i, P_j)$, with $i \ne j$, such that the line $P_iP_j$ contains a point $Q$ having integer coordinates and lying strictly between $P_i$ and $P_j$.

\item Let $a_1, a_2, \ldots, a_n, b_1, b_2, \ldots, b_n$ be $2n$ real numbers, where $a_1, a_2, \ldots, a_n$ are distinct, and suppose that there exists a real number $\alpha$ such that the product
\[(a_i + b_1)(a_i + b_2)\ldots(a_i + b_n)\]
has the value $\alpha$ for $i = 1, 2, \ldots, n$. Prove that there exists a real number $\beta$ such that the product
\[(a_1 + b_j)(a_2 + b_j)\ldots(a_n + b_j)\]
has the value $\beta$ for $j = 1, 2, \ldots, n$.

\item For nonnegative integers $n, r$, the binomial coefficient $\binom{n}{r}$ denotes the number of combinations of $n$ objects chosen $r$ at a time, with the convention that $\binom{n}{0} = 1$ and $\binom{n}{r} = 0$ if $n < r$. Prove the identity
\[\sum_{d=1}^{\infty} \binom{n-r+1}{d}\binom{r-1}{d-1} = \binom{n}{r}\]
for all integers $n$ and $r$, with $1 \le r \le n$.

\item Let $x$ be a real number with $0 < x < \pi$. Prove that, for all natural numbers $n$, the sum
\[\sin x + \frac{\sin 3x}{3} + \frac{\sin 5x}{5} + \ldots + \frac{\sin(2n-1)x}{2n-1}\]
is positive.

\item \begin{enumerate}
\item[(a)] The rectangle $PQRS$ has $|PQ| = \ell$ and $|QR| = m$, where $\ell$, $m$ are positive integers. It is divided up into $\ell m$ $1\times 1$ squares by drawing lines parallel to $PQ$ and $QR$. Prove that the diagonal $PR$ intersects $\ell + m - d$ of these squares, where $d$ is the greatest common divisor, $(\ell, m)$, of $\ell$ and $m$.
\item[(b)] A cuboid (or box) with edges of lengths $\ell$, $m$, $n$, where $\ell$, $m$, $n$ are positive integers, is divided into $\ell mn$ $1 \times 1 \times 1$ cubes by planes parallel to its faces. Consider a diagonal joining a vertex of the cuboid to the vertex furthest away from it. How many of the cubes does this diagonal intersect?
\end{enumerate}
\end{enumerate}
\end{document}
